\keepXColumns{}
\begin{tabularx}{\textwidth}{
  >{\raggedright\arraybackslash}p{1cm} 
  >{\raggedright\arraybackslash}p{2cm} 
  X
}
  \caption{Identifikasi Masalah Pengguna Perangkat Lunak Audit Energi}\label{tbl:problem} \\
  \toprule
  \textbf{Kode} & \textbf{Masalah} & \textbf{Deskripsi} \\
  \midrule
  \endfirsthead

  \caption{Identifikasi Masalah Pengguna Perangkat Lunak Audit Energi (lanjutan)} \\
  \toprule
  \textbf{Kode} & \textbf{Masalah} & \textbf{Deskripsi} \\
  \midrule
  \endhead

  \midrule
  \multicolumn{3}{l}{\textit{Bersambung ke halaman berikutnya}} \\
  %\bottomrule
  \endfoot

  \bottomrule
  \endlastfoot

  % === Table Data ===
  \hline
  P01 & Format data tidak seragam &
  Data berasal dari tagihan, \textit{spreadsheet}, dan pengukuran sementara dengan struktur berbeda, sehingga membutuhkan waktu untuk normalisasi dan validasi sebelum analisis \cite{luciano2011integrating,hamdani2017audit}. \\
  \hline
  P02 & Penginputan data manual &
  Penginputan data dimasukkan secara manual ke \textit{spreadsheet} sehingga meningkatkan risiko kesalahan ketik dan inkonsistensi antar auditor \cite{hamdani2017audit,winarto2019audit}. \\
  \hline
  P03 & Perhitungan tidak terdokumentasi &
  Formula dan metode perhitungan tidak baku atau terdokumentasi, menyebabkan hasil perhitungan berbeda antar sesi audit; otomasi perhitungan dan dokumentasi provenance dapat meningkatkan konsistensi dan keterjelasan hasil \cite{fitri2022sistem,kale2023provenance}. \\
  \hline
  P04 & Kurangnya histori terpusat &
  Hasil audit dan temuan disimpan secara terpisah atau dalam \textit{file} lokal sehingga sulit menelusuri riwayat dan menganalisis tren jangka panjang, yang menuntut sistem penyimpanan terpusat dan model data terstruktur \cite{hermawan2022perencanaan,hamdani2017audit}. \\
  \hline
  P05 & Pelaporan memakan waktu & Penyusunan laporan dilakukan secara manual sehingga proses ringkasan dan rekomendasi membutuhkan waktu lama dan menghambat tindak lanjut; otomatisasi laporan dapat mempercepat alur kerja \cite{putra2025evaluasi,winarto2019audit}. \\
  \hline
  P06 & Minim visualisasi dan \textit{insight} & Tidak tersedia \textit{dashboard} atau visualisasi terintegrasi sehingga pengguna kesulitan melihat pola konsumsi, beban puncak, dan area prioritas intervensi, sehingga diperlukan modul visualisasi yang informatif \cite{hamdani2017audit,hermawan2022perencanaan}. \\
  \hline
\end{tabularx}